% ============================================================
%  VORLAGE – Neue Initiative eintragen
%  Kopiere diese Datei und benenne sie um, z. B.:
%  cp _VORLAGE.tex meine_initiative.tex
%
%  Dann in main.tex eintragen:
%  \input{initiativen/meine_initiative}
% ============================================================

\begin{initiativkarte}{Name der Initiative}   % <── Initiativname hier

  % Bild (optional) – Datei in bilder/ ablegen
  \IfFileExists{bilder/meine_initiative.jpg}{%
    \begin{center}
      \includegraphics[width=\linewidth, height=5cm, keepaspectratio]{bilder/meine_initiative.jpg}
    \end{center}
    \vspace{0.4cm}
  }{}

  % Kurzbeschreibung
  \textit{\textcolor{hawgray}{Kurze Beschreibung in einem oder zwei Sätzen, die das Wesentliche der Initiative auf den Punkt bringt.}}

  \vspace{0.5cm}
  \hawrule

  % Ausführliche Beschreibung
  \textbf{\textcolor{hawblue}{Was machen wir?}}\\[0.3cm]
  Hier kommt die ausführlichere Beschreibung. Was sind die Ziele? Was macht ihr konkret? Welche Projekte habt ihr? Warum lohnt es sich beizutreten?

  \vspace{0.5cm}

  % Schlagwörter / Tags
  \tag{Schlagwort 1} \tag{Schlagwort 2} \tag{Schlagwort 3}

  \vspace{0.6cm}
  \hawrule

  % Info-Zeilen
  \infozeile{\faUsers}{Zielgruppe: Alle Studierenden, besonders Informatik, Ingenieurwesen}
  \infozeile{\faGraduationCap}{Semester: Ab 1. Semester}
  \infozeile{\faCalendarAlt}{Treffen: Jeden Dienstag, 18:00 Uhr, Raum XXX}
  \infozeile{\faEnvelope}{Kontakt: initiative@haw-hamburg.de}
  \infozeile{\faGlobe}{Website: \href{https://example.com}{example.com}}

\end{initiativkarte}
